% Document class and packages
\documentclass[a4paper, 12pt]{report}
\usepackage{geometry}
\usepackage[utf8]{inputenc}
\usepackage{graphicx}
\usepackage{titlesec}
\usepackage{fontenc}
\usepackage{float}
\usepackage{minted}
\usepackage[swedish]{babel}
\usepackage{pythonhighlight}
\usepackage{wrapfig}
% For headers and footers
\usepackage{fancyhdr, lastpage}
\usepackage{fancyvrb}
% Math spesific packages
\usepackage{amsfonts}
\usepackage{amsmath}
\usepackage{amssymb}
% For reffereces displaying and coloring as well as for some document settings
\usepackage[]{hyperref}
% For math symbols and figures
\usepackage{tikz}
% For color boxes
\usepackage{xcolor}
\usepackage[most]{tcolorbox}
% For unit displaying
\usepackage{siunitx}
% For correct paragraph phrasing
\usepackage{parskip}

% For figure caption formatting
\usepackage[font=small,labelfont=bf,center]{caption}

% For correct fotnotes and citation structure
\usepackage[backend=bibtex8,style=verbose,citestyle=alphabetic,bibstyle=alphabetic]{biblatex}
\addbibresource{ref.bib}

% Background image of title page
\usepackage{eso-pic}
\newcommand\BackgroundIm{
\put(0,0){
\parbox[b][\paperheight]{\paperwidth}{
\vfill
\centering
\includegraphics[height=\paperheight,width=\paperwidth,
keepaspectratio]{background-1.png}
\vfill
}}}

% To fix footnote counter between files
\usepackage{chngcntr}
\counterwithout*{footnote}{chapter}

% Global options
\geometry{a4paper}
\graphicspath{{Images/}}
\hypersetup{
  linkcolor=black,
  citecolor={blue!50!black},
  filecolor={blue!50!black},
  urlcolor={blue!80!black},
  pdftitle={The AES encryption algorithm},
  pdfauthor={Gabriel Lindeblad},
  pdfcreator={Gabriel Lindeblad},
  pdfsubject={Advanced Encryption Standard},
  pdfstartpage=1,
  pdffitwindow=true,
  colorlinks=true,
  pdfpagemode=UseNone,
  }
\setcounter{tocdepth}{3}
\setcounter{secnumdepth}{3}
\pagestyle{fancy}
\rhead{\currentchapter}
\cfoot{\thepage}
\fancyheadoffset[L, RO]{0cm}


% New commands and changed commands
\newcommand{\currentchapter}{}
\newcommand{\mychapter}[1]{\chapter{#1}
\renewcommand{\currentchapter}{#1}}
\renewcommand{\bibname}{Källförteckning}
\renewcommand{\contentsname}{Inehållsförteckning}
\addto\captionsswedish{%
  \renewcommand\listfigurename{Figur lista}}
\newcommand{\fakechapter}[2]{%
  \par\setcounter{chapter}{#2}
  \chaptermark{#1}
  \rhead{ }
}
\newcommand*{\captionsource}[2]{%
  \caption[{#1}]{%
    #1%
    \\\hspace{\linewidth}%
    \textbf{Källa:} #2%
  }%
}

% Titel reformating
\titleformat{\chapter}[block]
  {\normalfont\huge\bfseries}{\thechapter}{0.5em}{\Huge}
\titlespacing*{\chapter}{0pt}{-25pt}{0pt}

\titleformat{\section}[block]
  {\normalfont\Large\bfseries}{\thesection}{0.5em}{\Large}
\titlespacing*{\chapter}{0pt}{-25pt}{0pt}

\titleformat{\subsection}[block]
  {\normalfont\large\bfseries}{\thesubsection}{0.5em}{\large}
\titlespacing*{\chapter}{0pt}{-25pt}{0pt}

\titleformat{\subsubsection}[block]
  {\normalfont\normalsize\bfseries}{\thesubsubsection}{0.5em}{\normalsize}
\titlespacing*{\chapter}{0pt}{-25pt}{0pt}


% Glossaries setup
\usepackage[nopostdot,nonumberlist,style=altlisthypergroup,acronym,toc]{glossaries}
\makeglossaries

\renewcommand*{\glossaryname}{Begreppförklaring}
\renewcommand*{\acronymname}{Akronymer}
\renewcommand{\glsnamefont}[1]{\textbf{#1}}

%%%
\newglossaryentry{rsa}
{
    name={RSA},
    description={Rivest-Shamir-Adleman (RSA) är en av det mest välkända krypteringsalgoritmerna
    och var en av det första algorithmerna som byggde på en asymetrisk kryptering.
    RSA bygger på multiplikation av stora primatal där primtalen är
    nycklarna.\cite{rsa-ref}}
 }

\newglossaryentry{caesar}
{
    name={Caesarchiffer},
    description={Caesarchiffer är ett \gls{substitutionsskiffer},
    vilket helt enkelt bygger på att man byter ut varje bokstav i
    medelandet med en annan. Ersättnings bokstaven bestäms genom
    att man hoppar ett visst antal hopp i alphabetet som exempelvis
    3 hopp, vilket då innebär att ifall man har bokstaven a då skulle
    den bli ett d istället.\cite{caesar}}
 }

\newglossaryentry{xor}
{
    name={XOR},
    symbol={\(\oplus\)},
    description={XOR är en logisk operation inom datorvetenskap som fungerar ungefär som + uttrycket,
    med den enda skillnaden att {\(1 \oplus 1 = 0\)}. Detta samt att xor är en \gls{binär} operation, vilket
    innebär att termerna bara kan vara 0 eller 1 och resultatet det samma. Utöver XOR finns även
    \gls{or}, \gls{not} och \gls{and} bland annat.\cite{xor}} % Maybe needs to be refined
 }

\newglossaryentry{python}
{
    name={Python},
    description={Python är ett högnivå programmerings språk som är skrivet i programmerings språket C.
    Det är skapat av Guido van Rossum och släpptes i Februari 1991.\cite{python}}
 }

 \newglossaryentry{keystream}
 {
     name={Nyckelström},
     description={En nyckelström är i kryptografin en ström av \gls{pseudoslump}
     karaktärer som kan kombineras med exempelvis ett medelande för att producera
     en skiffertext.\cite{keystream}}
  }

  \newglossaryentry{pseudoslump}
  {
      name={Pseudoslump},
      description={Pseudoslump är en rad av nummer som kan se ut att vara helt
      slumpmässiga men har blivit framställda genom en upprepbar process.\cite{pseudoslump}}
   }

  \newglossaryentry{vscode}
  {
      name={VSCode},
      description={Visual Studio Code är en programutveklingsmiljö som är skapad av Microsoft.
      Det är ett öppet källkods projekt som är tillgängligt för det flesta operativsystem
      och kan användas för att skriva kod i flera olika språk.\cite{vscode-ref}}
   }

  \newglossaryentry{streamcipher}
  {
      name={Strömskiffer},
      description={Strömskiffer, ett symetriskt nyckel skiffer där man använder en
      \gls{pseudoslump}mässig skiffer ström (\gls{keystream}) som sedan en \gls{bit} i taget
      kombineras med det som ska krypteras. Den kombinerande operationen som används i
      strömskiffer är ofta en \gls{xor}-operation.\cite{streamcipher-ref}}
   }

  \newglossaryentry{bit}
  {
       name={Bit},
       description={En Bit är den minsta enheten av information som kan lagras i en dator.
       En bit kan endast ha två värden där den antingen är 0 eller 1, alltså ett \gls{binär}t
       värde. I datorvetenskap pratar man dock mer vanligen om ett \gls{byte} som är 8 bits.\cite{bit-ref}}
   }

  \newglossaryentry{byte}
  {
      name={Byte},
      description={En Byte består av 8 \gls{bit}s och är en enhet som används inom
      datorvetenskap. En byte kan ha 256 olika värden från 0 till 255, vilket är $2^8$ värden.
      Dessa värden representerar ofta tecken eller symboler som exempelvis bokstäver, siffror med mera.
      Tolkningen av vad en sekvens av bytes eller en enskild byte står för beror däremot på vilken
      teckenkodning som används. Exempel på teckenkodningar kan var \acrshort{ascii} och ISO-8859.\cite{byte-ref}}
   }

  \newglossaryentry{hashfunktion}
  {
      name={Hashfunktion},
      description={Hashfunktion är en funktion som man kan säga delar upp en viss datamängd
      och genom för en serie operationer som resulterar i hashtext av godkännd längd.
      Längden är samma för alla hastexter som använder samma funktion medans innehållet
      förändras så fort något över huvud taget ändras i datan som funktionen appliceras på.
      Användningsområdet för dessa funktioner är bland annat när man vill kunna verifiera
      medelanden eller information och försäkra sig om att ingen ändrat på medellandet efter att det skickats.
      Detta kan man då göra för att man vet att om man skör informationen genom samma hashfunktion
      borde resultatet vara identisk ifall informationen är oförändrad.\cite{hashfunktion-ref}}
   }

  \newglossaryentry{polyalphabetic-substitutionsskiffer}
  {
      name={Polyalphabetic substitutionsskiffer},
      description={polyalphabetic substitutionsskiffer bygger på att man använder flera olika \gls{substitutionsskiffer} för att
      på så sätt undvika en ut av det största svagheterna med \gls{substitutionsskiffer}. Detta då att dom lätt går att
      knäcka genom en \gls{frekvensanalys} då vissa bokstäver dyker upp mer frekvent i språket än andra. För att lösa
      detta så använder polyalphabetiska skiffer flera olika substitutionsskiffer som man byter mellan med en viss
      frekvens för att eliminera \gls{frekvensanalys}ens effektivitet.\cite{polyalphabetic-ref}}
   }

  \newglossaryentry{frekvensanalys}
  {
    name={Frekvensanalys},
    description={Frekvensanalys inom kryptografi är en metod för att knäcka ett \gls{substitutionsskiffer} genom att
    analysera frekvensen av bokstäver och utnyttja det faktum att en del bokstäver framkommer mer frekvent än andra
    i språket. På detta viset kan man då sedan lista ut vilka bokstäver som är vilka i det krypterade medellandet.\cite{frekvensanalys-ref}}
    }

  \newglossaryentry{substitutionsskiffer}
  {
      name={Substitutionsskiffer}, % fix this
      description={Ett Substitutionsskiffer är en typ av krypterings metod som bygger på att man
      byter ut delar av informationen man ska kryptera med exempelvis andra symboler med hjälp av en nyckel.
      Detta kan exempelvis vara bokstäver som byts ut mot andra bokstäver precis som i \gls{caesar} eller
      siffror som byts ut mot andra siffror.\cite{substitutionsskiffer-ref}}
    }

  \newglossaryentry{and}
  {
      name={AND},
      description={AND är en logisk operation inom datorvetenskap och matematik
      som tar två \gls{binär}a värden och ger till baka ett \gls{binär}t värd. Detta värde
      är 1 om och endast om båda värdena är 1 annars är värdet 0.\cite{logical-operators-ref}}
    }

  \newglossaryentry{or}
  {
      name={OR},
      description={OR är en logisk operation inom datorvetenskap och matematik
      som tar två \gls{binär}a värden och ger till baka ett \gls{binär}t värd. Detta värde
      är 1 om minst ett av värdena är 1 annars är värdet 0.\cite{logical-operators-ref}}
    }

  \newglossaryentry{not}
  {
      name={NOT},
      description={NOT är en logisk operation inom datorvetenskap och matematik
      som tar två \gls{binär}a värden och jämför dom. Det slutgiltiga värdet som ges tillbaka är
      1 om värdena inte är lika varandra och annars är värdet 0.\cite{logical-operators-ref}}
    }

  \newglossaryentry{binär}
  {
      name={Binär},
      description={Binär är ett begrepp som används inom datorvetenskap och matematik för att beskriva
      ett värde som endast kan vara två olika saker som exempelvis 0 eller 1, sant eller falskt. Detta innebär abstract
      binära tal bygger på talbasen 2, vilket skiljer sig från vad som ofta vanligen används i vardagen
      som är talbasen 10.\cite{binär-ref}}
    }

  \newglossaryentry{enigma}
  {
      name={Enigma},
      description={Enigma var ett krypterings verktyg som användes under andra världskriget av tyska
      militären för att kryptera medelanden. Maskinen bestod av en elektromekanisk rotordisk som under tiden
      medellandet skrivs in för kryptering även ändrar det elektriska kopplingarna mellan vilka bokstäver som blir vad.
      Detta är en av sakerna som gjorde enigma väldigt svår att knäcka samt en av anledningarna till att
      liknande maskiner användes under stora delar av det tidiga 1900-talet.\cite{enigma-ref}}
    }
  \newglossaryentry{SP-network}
  {
      name={SP-network},
      description={SP-network eller även kallat Substitution-permutation network är inom kryptografin
      en serie av matematiska operationer som genomförs i rundor för att på så sätt kryptera ett medelande. Det består
      ut av två delar, en substitutions del och en permutation del. Substitutionen delen fungerar
      precis som \nameref{sec:aes-subbytes} medans permutationen exempelvis skulle kunna representeras
      med \nameref{sec:aes-shiftrows}.\cite{sp-network-ref}}
    }
  \newglossaryentry{ssh}
  {
      name={SSH},
      description={SSH eller som det även är kallat Secure Shell är ett protokoll som används för säker kommunikation
      över ett nätverk. SSH är ett exempel på ett nätverksprotokoll och används bland annat för att kryptera medelanden
      som skickas mellan två enheter.\cite{ssh}}
    }
  \newglossaryentry{http}
  {
        name={HTTP},
        description={HTTP (Hypertext Transfer Protocol) är det protokoll som används när man besöker en webbsida.
        Detta protokoll har används sedan 1990-talet och används fortfarande till stora delar.\cite{http}}
    }
  \newglossaryentry{hexadecimal}
  {
        name={Hexadecimal},
        description={Hexadecimal tal är tal med talbasen 16. Detta innebär att det finns 16 olika symboler som kan användas
        istället för 10 symboler som används i decimaltal som har talbasen 10. Detta görs genom att man representerar tal
        större än 9 med bokstäverna A-F.\cite{NE-2}}
    }
  \newglossaryentry{ppm}
  {
        name={PPM},
        description={PPM eller även kallad Portable Pixel Map är ett fil format som förvara bilder i råa bytes som då representerar
        färgerna i bilden. PPM använder sig av 3 \gls{byte}s för att representera en färg.\cite{ppm}}
    }
  \newglossaryentry{gimp}
  {
        name={GIMP},
        description={GIMP är ett open source bild redigerings program som är tillgängligt för det flesta operativsystem.\cite{gimp}}
    }
  \newglossaryentry{win11}
  {
        name={Windows 11},
        description={Windows 11 är den senaste versionen av ett \gls{operativsystem} för datorer framtaget av företaget Microsoft. Windows 11
        tillhör \gls{operativsystem} familjen Windows NT som först lanserades 1993.\cite{win11}}
    }
  \newglossaryentry{operativsystem}
  {
        name={Operativsystem},
        description={Operativsystem är det program på en dator som fungerar som gränssnitt mellan användaren och datorns maskinvara. Detta
        innebär att operativsystemet hanterar saker som inmatning från användaren till att tillhandahålla ett gränssnitt för olika program
        så att dom kan genomföra sina uppgifter.\cite{opsystem}}
    }
  \newglossaryentry{cpuh}
  {
      name={CPU klockhastighet},
      description={CPU klockhastighet är den hastighet som en CPU kan utföra instruktioner på. Den mäts i hertz och är frekvensen som
      klockgeneratorn i en dators CPU använder sig av för att bland annat styra hur snabbt instruktioner ska utföras samt synka
      olika komponenter i datorn.\cite{cpuh}}
    }
  \newglossaryentry{matrismultiplikation}
  {
      name={Matrismultiplikation},
      description={Matrismultiplikation är en matematisk metod för att multiplicera ihop två eller flera matriser med varandra. Detta görs
      genom att värdena i den första matrisen multipliceras med värdena i den andra matrisen och sedan summeras ihop. Detta resulterar
      i en ny matris som är resultatet av multiplikationen.\cite{matrismultiplikation}}
    }

\newacronym{aes}{AES}{Advanced Encryption Standard}
\newacronym{ecb}{ECB}{Electronic Code Book läge}
\newacronym{cbc}{CBC}{Cipher Block Chaining läge}
\newacronym{ofb}{OFB}{Output Feedback läge}
\newacronym{des}{DES}{Data Encryption Standard}
\newacronym{iv}{IV}{Initialization Vector}
\newacronym{ssl}{SSL}{Secure Socket Layer}
\newacronym{tls}{TLS}{Transport Layer Security}
\newacronym{ascii}{ASCII}{American Standard Code for Information Interchange}
\newacronym{wpa2}{WPA2}{Wi-Fi Protected Access 2}
\newacronym{nist}{NIST}{National Institute of Standards and Technology}
\newacronym{rcon}{Rcon}{Round Constant}
%%%

% Data visualization setup
\usepackage{pgfplots}
\usetikzlibrary{arrows.meta,decorations.pathmorphing,backgrounds,positioning,fit,petri}
%%%

% Start document
\begin{document}

\newgeometry{top=0cm, bottom=2cm, left = 1.5cm, right = 1.5cm}

% Beginning title page
\begin{titlepage}
    % Inserting background images
    \AddToShipoutPicture*{\BackgroundIm}

    % School logo
    \begin{figure}[H]
    \includegraphics[width=0.3\textwidth]{skola_logga.png}
    \end{figure}

    % Aligning year to the right
    \raggedleft

    % Year
    \vspace{-3.5cm}
    {\large \textbf{2022/2023}}

    % Aligning to the left
    \raggedright

    % Title and sub title
    \vspace{3.5cm}
	{\huge\bfseries The AES encryption algorithm\par}
   ~{\large En undersökning av The \acrfull{aes}}

% Bottom of the page
    \vfill

  {Klass:} \\
    ~\textbf{NA20}

  {Handledare:} \\
    ~\textbf{Jimmy Nylén}

  {Författare:} \\
    ~\textbf{Gabriel Lindeblad}

  {Program:} \\
    ~\textbf{Naturvetenskapsprogrammet}

    \vspace{1cm}
	{\large \today\par}
    \vspace{-1cm}

% End of title page
\end{titlepage}

% Geometry of page borders for abstract
\newgeometry{top=5cm, bottom=2cm, left = 2cm, right = 2cm}

% Abstract
\chapter*{Abstract}
This paper is a study of the \acrlong{aes} with the purpose to investigate how the key length and running mode affects the encryption speed of the algorithm
as well as how the different running modes affect the security of the algorithm. The study was conducted by implementing \acrshort{aes} and the running modes \acrshort{ecb}, \acrshort{cbc} and \acrshort{ofb} in the programming language \gls{python} and
then preforming a series of tests on it. These tests consisted of measuring the encryption speed of the algorithm with different key lengths and running modes as well as the encryption
of a picture file using each of the running modes. The result of the study showed that the encryption speed of the algorithm increased as the key length increased, while a noticeable speed difference
between the different running modes not could be assessed from the result of the study. The result of the study also showed that the running mode \acrshort{ecb} was the least secure of the three
running modes tested.

% New page borders for bulk of the report
\newgeometry{top=2.5cm, bottom=2cm, left = 2cm, right = 2cm} % 2.25cm is the height of the header (is a checky solution that might need to bee removed if more headings are added)

% Inehållsförteckning
\tableofcontents

\listoffigures

\clearpage

\printglossary

\printglossary[type=\acronymtype, style=super]

\mychapter{Inledning} % Mostly clear but might need some refinement
Kryptering, en bärande grundsten i dagens digitaliserade samhälle. Det är väggen mellan oss och resten
av världen, ett lås runt våra liv. Kryptering bygger på ett simpelt koncept, att dölja informationen
från all förutom den menade mottagaren. Ett koncept som har funnits med oss under stora delar av
människans historia och som än idag är en viktig del av vårt samhälle.\footfullcite{luciano1987cryptology}

Att dölja information har gått från de enkla metoder som användes redan för 2000 år sedan för att förmedla hemlig information med hjälp av
exempelvis \gls{caesar} till det komplexa och invecklade algoritmer som idag används för att skydda nästan all information som lagras och skickas
över internet.\footcite{luciano1987cryptology} Dessa komplexa algoritmer har utvecklats under en tid där datorer står som det dominerande informationshanteringsverktyget och
därför är det även en dator som används som huvudverktyg i denna rapport för att undersöka en av det vanligaste krypteringsalgoritmerna i dagens samhälle,
The \acrlong{aes}.

\section{Syfte} % Mostly clear but might need some refinement
Syftet med undersökning är att undersöka krypterings algoritmen \acrshort{aes},
för att utveckla en fördjupad förståelse för mer avancerade krypterings algoritmer.
Samt bygga en uppfattning om hur det på olika sätt går att implementera krypterings algoritmer
och vilken påverkan detta då får på deras prestanda och säkerhet.

\section{Frågeställningar} % Mostly clear but might need some refinement and reformulating
\begin{itemize}
    \setlength{\itemindent}{-1em}
    \item Hur påverkas tiden det tar att kryptera något mellan det olika nyckel längderna 128-bit,
          192-bit och 256-bit nyckel?

    \item Hur förändras krypterings tiden mellan de olika körlägena \acrshort{ecb}, \acrshort{cbc} \& \acrshort{ofb}?

    \item Hur påverkas skiffertexten av det olika körlägena \acrshort{ecb}, \acrshort{cbc} \& \acrshort{ofb}
          samt vilken betydelse får det ur ett säkerhetsperspektiv?
\end{itemize}

\section{Avgränsning} % Mostly clear but might need some refinement and reformulating

Denna undersökning är avgränsad till att endast undersöka \acrshort{aes} och dess
användning med fokus på hur nyckellängd och körläge påverkar krypteringstiden. Detta samt hur
den resulterande skiffer texten påverkas av det olika körlägena \acrshort{ecb}, \acrshort{cbc} \& \acrshort{ofb}
ur ett säkerhetsperspektiv.

Denna analys av algoritmens säkerhet tar inte hänsyn till faktorer så som möjliga attacker där
ibland exempelvis Brute-Force\footfullcite{kumar2011investigations}
\& Side-Channel\footfullcite{10.1007/978-3-642-04138-9_8} attacker. Undersökningen är även
begränsad till att endast utföras på en mjukvaruimplementering av \acrshort{aes}.

\chapter{Bakgrund}

\section{section 1}
abstract
\section{section 2}
abstract
\section{section 3}
abstract

\mychapter{Teori}


\section{Kryptering}


\section{Blockchiffer}


\subsection{Körlägen}


\subsubsection{ECB}
\acrfull{ecb} är en av det enklaste blockchiffer körlägena som finns.
\acrshort{ecb} i sig är ganska lätt att förstå och bygger i huvudsak bara på
att man delar upp den data man vill kryptera i delar kallade block och tar sedan varje
block för sig och kör genom algoritmen, vilket tydligt visas i
Figur \ref{fig:ecb-mode-enc} \& \ref{fig:ecb-mode-dec}.
\footfullcite{modesofoperation}

Figur \ref{fig:ecb-mode-enc} visar hur \acrshort{ecb} fungerar vid kryptering.
Här visas hur varje block för sig krypteras med hjälp av en blockchiffer algoritm
tillsammans med den givna nyckeln.

\begin{figure}[H]
    \includegraphics[width=\textwidth]{ECB_encryption.png}
    \captionsource{\acrlong{ecb} kryptering}{{\cite{ecb-mode-enc-ref}}}
    \label{fig:ecb-mode-enc}
\end{figure}

Figur \ref{fig:ecb-mode-dec} visar istället hur \acrshort{ecb} fungerar vid
dekryptering, vilken är en till stort sett identisk operation med det enda undantaget
att blockchiffret körs i dekrypterings läge istället för krypterings läge.

\begin{figure}[H]
    \includegraphics[width=\textwidth]{ECB_decryption.png}
    \captionsource{\acrlong{ecb} dekryptering}{{\cite{ecb-mode-dec-ref}}}
    \label{fig:ecb-mode-dec}
\end{figure}

På grund av \acrshort{ecb} körlägets simplicitet så finns det dock även ett ganska
stort problem med detta körläge. Det handlar om att \acrshort{ecb} inte på något
sätt förhindrar att två block med samma innehåll som krypteras inte resulterar i
ett identiskt krypterat block.\footcite{modesofoperation}

Vad detta innebär är att för större mängder data
är att det börjar bildas mönster i skiffertexten. Detta är något som väldigt
tydligt visar sig ifall man krypterar en bild, vilket går att se när man jämför
figur \ref{fig:pi-original} \& \ref{fig:pi-ecb}.
Det här faktumet är även varför \acrshort{ecb} inte är ett säkert körläge
och därför inte används näst intill aldrig i praktiken.\footcite{modesofoperation}

\subsubsection{CBC}
\acrlong{cbc} är ett av de mest vanligen använda körlägena för många blockchiffer.

\footcite{modesofoperation}

\begin{figure}[H]
    \includegraphics[width=\textwidth]{CBC_encryption.png}
    \captionsource{\acrlong{cbc} kryptering}{{\cite{cbc-mode-enc-ref}}}
    \label{fig:cbc-mode-enc}
\end{figure}

\begin{figure}[H]
    \includegraphics[width=\textwidth]{CBC_decryption.png}
    \captionsource{\acrlong{cbc} dekryptering}{{\cite{cbc-mode-dec-ref}}}
    \label{fig:cbc-mode-dec}
\end{figure}

\subsubsection{OFB}


\section{Symetrisk \& Asymmetrisk Kryptering}
Symetrisk och asymetrisk kryptering handlar om hur nycklar används i olika
krypteringsalgoritmer. För symetriska krypterings algoritmer så betyder detta
att samma nyckel är vad som används för både kryptering och dekryptering.

\footfullcite{symencrypt}

\section{AES}


\subsection{Finite Fields}


\subsection{AES S-Box}


\subsection{Struktur}


\subsubsection{SubBytes operation}
SubBytes operationen bygger på ...

\subsubsection{ShiftRows operation}


\subsubsection{MixColumns operation}


\subsubsection{AddRoundKey operation}


\subsection{Nyckel utökning}


\subsubsection{RotWord}


\subsubsection{SubWord}


\subsubsection{Rcon}


\subsection{AES-128bit}


\subsection{AES-192bit}


\subsection{AES-256bit}



\mychapter{Metod \& Genomförande} % (To be reviwed)
Metoden för denna undersökning bygger på en implementering av \acrshort{aes} i programmeringsspråket
\gls{python}. Detta tillsammans med ett antal konstruerade tester även dom implementerade i
\gls{python} är vad som använts för själva undersökningen av \acrshort{aes}. Koden
är skriven med hjälp av programmet \gls{vscode} och är byggd huvudsakligen för \gls{python} 3.10 men
på grund av att \gls{python} 3.11 släpptes innan undersökningen genomfördes så är de istället \gls{python} 3.11.0
som användes under undersöknings genomförandet.\footfullcite{python311}

\section{Implementering} % (To be reviwed)
Implementeringen av \acrshort{aes} är uppdelad i ett antal funktioner till stor del är baserat på
hur strukturen och uppdelningen av \acrshort{aes} beskrivs i \citetitle{daemen1999aes}\footcite{daemen1999aes}.
Några av det Huvudsakliga funktionerna av algoritmen är som följande:
\begin{itemize}
    \item \texttt{\nameref{sec:aes-subbytes}}
    \item \texttt{\nameref{sec:aes-shiftrows}}
    \item \texttt{\nameref{sec:aes-mixcolumns}}
\end{itemize}

Dessa funktioner används i varje runda och utgör grunden av algoritmen. Utöver detta finns
även \nameref{sec:aes-addroundkey} som används för att lägga till nyckeln i varje runda. AddRoundKey steget är dock inte
en funktion eftersom det enbart består ut av en \gls{xor}-operation. Utöver detta finns det även
en funktion som används för att expandera nyckeln som beskrivs i \nameref{sec:aes-key-expansion} samt
individuella funktioner för kryptering och dekryptering till de olika körlägen \nameref{sec:ecb}, \nameref{sec:cbc} och
\nameref{sec:ofb}.

Implementeringen använder sig av \gls{python} biblioteket NumPy för att hantera matriserna av \gls{byte}s och genomföra
matematiska operationer på dessa. Hela implementeringen går att hitta i appendix \ref{app:python} "\nameref{app:python}" samt
förklaras även i avsnittet \nameref{sec:aes}.

\section{Test Uppsättning} % (To be reviwed)
Test uppsättningen är uppdelad i det tre delarna \nameref{sec:nyckellangd-test}, \nameref{sec:körlages-test} och \nameref{sec:krypterings-test}. Två av delarna, \nameref{sec:nyckellangd-test} och
\nameref{sec:körlages-test} är skrivna som ett \gls{python} script vilket visas i appendix \ref{app:analyze}. När det gäller Krypteringstest så skiljer sig det från de två andra analyserna eftersom det
inte är en automatiserad process utan ett antal praktiska steg som genomfördes för att generera resultatet.

\subsection{Testmiljö} % Describe the test enviroment (might reqiuer extra glossary entrys for some of the terms)
Genomförandet av undersökningen skedde på en Windows dator med \gls{python} 3.11.0 installerat. I syfte att minimerar variabiliteten
mellan de olika omgångarna av tester då allt repeteras ett vist antal gånger så begränsades vissa delar av datorn och anpassningar gjordes i genomförandet.
De begränsningar som placerades på datorn bestod ut av att datorn CPU klockhastighet begränsades till sin grund hastighet för att ge
varje test ett likvärdigt utgångsläge som möjligt. Utöver detta så genomfördes även varje test separat utan parallellisering för att
eliminera risken av att testerna skulle kunna påverka varandra.

\subsection{Nyckellängds Test} % Explain the key length test
\label{sec:nyckellangd-test}
Nyckellängds testet bygger på att varje nyckellängd används för att kryptera en text fil på 1 MB. Detta genomförs 25 gånger för
varje nyckellängd och resultatet av varje test sparas sedan ...

\subsection{Körläges Test} % Explain the mode of operation test
\label{sec:körlages-test}


\subsection{Krypterings Test} % explain how the encryption test is setup
\label{sec:krypterings-test}
Krypterigns testet består utav att man krypterar en bild i bild formatet \gls{ppm} med hjälp av \acrshort{aes} i tre olika körlägen.
PPM formatet används då de tack vare sin enkla struktur kan genomgå en kryptering och fortfarande representeras som en bild
efteråt. Detta gör det då möjligt att visualisera resultatet av krypteringen. För att kryptera bilden används den implementerade
\acrshort{aes} algorithmen som körs i körlägena \nameref{sec:ecb}, \nameref{sec:cbc} och \nameref{sec:ofb} körlägen.

Före och efter krypteringen modifieras dock \gls{ppm} filen så att den går att öppna när den väl är krypterad. Detta görs genom att de fyra
första raderna i filen tas bort och sätts sedan tillbaka efter krypteringen. Detta gör att filen fortfarande kan öppnas som en bild
efter krypteringen.

\section{Genomförande} % Describe how the test was preformed (To be reviwed)
Först kördes Nyckellängdstestet och Körlägestestet genom att köra \gls{python} filen analyze.py. Där efter
genomfördes Krypterings testet. För krypteringstestet valdes först en lämplig bild ut för krypteringen. Sedan
konverterades bilden till fil formatet \gls{ppm} med hjälp av programmet \gls{gimp}. Där efter utförs följande kommandon i en terminal:

\begin{minted}{bash}
    $ head -n 4 pi.ppm > header.txt

    $ tail -n +5 pi.ppm > body.bin

    $ python3 encrypt.py 2b7e151628aed2a6abf7158809cf4f3c body.bin ECB

    $ cat header.txt body.bin.enc > pi-ecb.ppm
\end{minted}

Här används file header.txt som en temporär förvaringsplats för de första fyra raderna i \gls{ppm} filen. Detta
efter som de första fyra raderna innehåller informationen som berättar hur filen senare ska tolkas. Där efter används
tail kommandot för att ta bort de första fyra raderna i filen och skriva resten till en ny fil med namnet body.bin.
Filen body.bin krypteras sedan med hjälp av \gls{python} filen encrypt.py med nyckeln 2b7e151628aed2a6abf7158809cf4f3c.

Innehållet från den nya filen body.bin.enc sätt sedan ihop med innehållet av header.txt som infogas som de fyra första raderna samt filen döps om till pi-ecb.ppm.
Slutligen så konverteras även filen tillbaka till ett mer lätthanterligt format med hjälp av \gls{gimp}.
Processen beskriven ovan upprepades sedan för de två andra körlägen \nameref{sec:cbc} och \nameref{sec:ofb} med
den enda skillnaden att de istället för ECB används CBC och OFB som argument till encrypt.py filen (Det tredje terminal kommandot).

Därefter påbörjades \nameref{sec:nyckellangd-test} och \nameref{sec:körlages-test}, vilket utfördes genom att köra
\gls{python} filen analyze.py. Filen analyze.py genomförde då krypteringen av en textfil på 1MB fylld med slumpmässigt genererade \gls{byte}s med hjälp av \acrshort{aes}
implementeringen som går att se i appendix \ref{app:python}. Detta genomfördes i de tre olika körlägena \acrshort{ecb}, \acrshort{cbc} och \acrshort{ofb} samt de tre olika nyckellängderna 128bit, 192bit och 256bit.
För varje nyckellängd och körläge så genomfördes krypteringen 25 gånger. Resultatet exporterades därefter till 6 olika text filer, vilkas innehåll sedan sammanställdes i \nameref{chap:resultat}.

\mychapter{Resultat}
\label{chap:resultat}
\section{Nyckellängds Test}
\label{sec:nyckellangd}

\section{Körläges Test}
\label{sec:korlages}

\section{Krypterings Test}
\label{sec:krypterings-test}


\mychapter{Diskussion} % Diskussion av resultat, presentera och tolka resultatet samt för ett nyanserat resonemang om vad resultatet betyder. Diskutera även eventuella brister i experimentet och hur dessa kan förbättras.
\label{chap:discussion}

% Tolka resultatet från nyckellängds undersökningen
Utifrån resultatet från undersökningen kan man då se hur tiden det tar att kryptera en fil på 1MB ökar när man använder längre nycklar så som 192-\gls{bit} 256-\gls{bit} nycklar
i förhållande till en 128-\gls{bit} nyckel. Något som bland annat visar sig när man tittar på den procentuella tidskillnaden mellan de olika nyckellängderna. Där man kan se en
tids skillnad på 18,8\% mellan 128-\gls{bit} \& 192-\gls{bit} nyckeln och 27,8\% mellan 128-\gls{bit} \& 256-\gls{bit} nyckeln.
Anledningen till detta är något som högst troligen beror på att antalet krypteringsrundor som genomförs för det olika nyckellängderna ökar när nyckeln blir längre.
Något som då höjer antalet operationer som genomförs för varje 16-\gls{byte} block som krypteras, vilket i sin tur höjer den totala krypterings tiden.

% Tolka resultatet från körläges undersökningen
Resultatet visar även att skillnaden i krypteringstid mellan olika körlägen är väldigt liten. Mellan \acrshort{ecb} \& \acrshort{cbc} var skillnaden 0,3\% och mellan \acrshort{ecb}
\& \acrshort{ofb} var skillnaden 0,4\%. Något som inte riktigt var väntat då det för både \acrshort{cbc} och \acrshort{ofb}
genomförs en ytterligare operation mellan varje 16-\gls{byte} block som krypteras för att länka ihop de olika blocken. Men samtidigt så skulle anledning till att skillnaderna är
så små kunna vara att \gls{xor}-operationen som används för att länka ihop blocken är en operation som är relativt lätt och snabb för en dator att genomföra, samtidigt som
operationen även bara behöver genomföras en gång för varje 16-\gls{byte} block som krypteras.

I resultatet av körläges testet kan man dock även se en viss spridning i tiden som det tar att kryptera en fil mellan det olika omgångarna.
Något som visar på en viss osäkerhet då \acrshort{ecb} har en högsta tid som är högre än både den högsta tiden för \acrshort{cbc} och \acrshort{ofb}. Vilket då skulle kunna antyda
att de skillnader som mätts upp mellan medelvärdena kan vara något som är inom felmarginalen för mätningen. Detta innebär då att det inte går att säga med säkerhet att
det finns någon större skillnad i krypterings tiden mellan det olika körlägena utifrån denna undersökning.

% Tolka resultatet från krypteringsmetod undersökningen
Tittar man sedan istället på hur säkerheten påverkas av det olika körlägena så kan man ganska tydligt se att \acrshort{ecb} är det mest osäkra körläget för större data mängder.
Detta då man som visas i resultatet kan se hur trots kryptering man fortfarande kan se spår av den ursprungliga bilden i informationen som krypterats. Något som inte
går att göra när bilden istället krypteras med hjälp av \acrshort{cbc} eller \acrshort{ofb}.

% Utvärdera metoden (nyanserat resonemang)



\section{Felkällor} % Variabilitet i klockhastighet, Ogämnheter i implementeringen som resulterar i missvisande tider, osäkerhet host mätmetoden, fel i implementeringen,


\section{Förbättringar} % Öka antalet reppetitioner för att minskar risken för fel, använda andra metoder för att mäta tiden, testa olika nycklar för att se hur de påverkar,


\section{Slutsats}



\fakechapter{Källförteckning}{7}
\addcontentsline{toc}{chapter}{Källförteckning}
\printbibliography[title=Källförteckning]

\clearpage

\appendix

\chapter{AES körlägestest resultat}
\label{app:raw-data-mode-picture-test}

\section{Före testet}
\begin{figure}[H]
  \centering
  \includegraphics[width=0.6\textwidth]{pi.png}
  \caption{Orginal bild}
  \label{fig:pi-original}
\end{figure}

\section{Efter ECB kryptering}
\begin{figure}[H]
  \centering
  \includegraphics[width=0.6\textwidth]{pi-ecb.png}
  \caption{Efter ECB Kryptering}
  \label{fig:pi-ecb}
\end{figure}

\section{Efter CBC kryptering}
\begin{figure}[H]
  \centering
  \includegraphics[width=0.6\textwidth]{pi-cbc.png}
  \caption{Efter CBC Kryptering}
  \label{fig:pi-cbc}
\end{figure}

\section{Efter OFB kryptering}
\begin{figure}[H]
  \centering
  \includegraphics[width=0.6\textwidth]{pi-OFB.png}
  \caption{Efter OFB Kryptering}
  \label{fig:pi-ofb}
\end{figure}

\cleardoublepage

\chapter{Fullständig data från testerna}
\label{app:raw-data}

\section{Data från nyckellängdstest}
\label{app:raw-data-keylength}

\begin{table}[H]
    \centering
    \begin{tabular}{ ||c|c|c|| }
      \hline
      \multicolumn{3}{|c|}{\bfseries{Resultat Tid (s)}} \\
      \hline
      \bfseries{ECB - 128bit} & \bfseries{ECB - 192bit} & \bfseries{ECB - 256bit} \\
      \hline
      21.114332500001183 & 25.019810000005236 & 29.25573820000136 \\
      21.2294222000055 & 25.03902250000101 & 29.199436000002606 \\
      21.400101300001552 & 24.921688300004462 & 29.1446350000042 \\
      21.394803399998636 & 25.010479399999895 & 29.276371100000688 \\
      21.410092899997835 & 25.023025200003758 & 29.208361999997578 \\
      21.255479499996 & 24.94711079999979 & 29.12308629999461 \\
      21.709606899996288 & 24.932161500000802 & 29.057567500000005 \\
      20.930029600000125 & 24.93615460000001 & 29.23549970000022 \\
      20.91164679999929 & 25.07268920000206 & 29.163480200004415 \\
      21.024735999999393 & 25.05104990000109 & 29.13999270000204 \\
      20.974247600002855 & 24.92792770000233 & 29.03159289999894 \\
      20.906690399999206 & 25.040346800000407 & 29.238432199999806 \\
      20.949946799999452 & 25.02671339999506 & 29.226332500002172 \\
      20.96879309999349 & 24.999294099994586 & 29.13317239999742 \\
      20.955687699999544 & 24.973698000001605 & 29.119369399995776 \\
      20.907472200000484 & 24.992430599995714 & 29.273320899999817 \\
      20.99763510000048 & 25.22444600000017 & 29.147454799996922 \\
      20.964201399998274 & 25.310352100001182 & 29.133893199999875 \\
      20.970198600000003 & 25.03863189999538 & 29.14858010000171 \\
      21.033452000003308 & 24.887886799995613 & 29.20511969999643 \\
      20.835713899999973 & 25.002411200002825 & 29.214967400002934 \\
      20.933525599997665 & 25.029954099998577 & 29.155635700000857 \\
      20.899700900001335 & 24.92795739999565 & 29.136839000006148 \\
      20.861958300003607 & 25.015061800004332 & 29.250083099999756 \\
      20.898621000000276 & 25.012236000002304 & 29.22063089999574 \\
      \hline
    \end{tabular}
  \end{table}

\section{Data från körlägestest}
\label{app:raw-data-mode}

\begin{table}[H]
    \centering
    \begin{tabular}{ ||c|c|c|| }
      \hline
      \multicolumn{3}{|c|}{\bfseries{Resultat Tid (s)}} \\
      \hline
      \bfseries{ECB} & \bfseries{CBC} & \bfseries{OFB} \\
      \hline
      20.855055499996524 & 20.863876399998844 & 20.801168300000427 \\
      20.88846609999746 & 20.907602499995846 & 20.78754040000058 \\
      20.879108299996005 & 20.86540999999852 & 20.823867599996447 \\
      20.801008999995247 & 20.896374500000093 & 20.87718119999772 \\
      20.85887470000307 & 20.868871400001808 & 20.770114799997828 \\
      20.87116899999819 & 20.882914800000435 & 20.767576799997187 \\
      20.854077200005122 & 20.89511720000155 & 20.85902029999852 \\
      20.85488130000158 & 20.85610540000198 & 20.797623000005842 \\
      20.826587200004724 & 20.865844100000686 & 20.796596600004705 \\
      20.909639099998458 & 20.85497480000049 & 20.797125700002653 \\
      20.870294899999863 & 20.906652400000894 & 20.835499799999525 \\
      20.818692400003783 & 20.92016809999768 & 20.865253799995116 \\
      20.82849349999742 & 20.908798900003603 & 20.800847200000135 \\
      20.86879109999427 & 20.850567699999374 & 20.795617399999173 \\
      20.91046159999678 & 20.883563499999582 & 20.834485100000165 \\
      20.82723100000294 & 20.884897599993565 & 20.82335030000104 \\
      20.80391700000473 & 20.89037549999921 & 20.828929399998742 \\
      20.87521669999842 & 20.885131600000022 & 20.811529400001746 \\
      20.91481080000085 & 20.93468269999721 & 20.881663000000117 \\
      21.085760299996764 & 20.859826599997177 & 20.85081160000118 \\
      20.97296980000101 & 20.875606900001003 & 20.823041400006332 \\
      21.053496200001973 & 20.903108300000895 & 20.791531299997587 \\
      20.978821899996547 & 20.892179500006023 & 20.878361099996255 \\
      20.954259099999035 & 20.890548900002614 & 20.873186500000884 \\
      21.028453000006266 & 20.88354160000017 & 20.76938550000341 \\
      \hline
    \end{tabular}
  \end{table}

\chapter{AES Implementering i python}
\label{app:python}

\section{AES.py}

\begin{python}
# ---------------
# Imported libraries
# ---------------
from os.path import getsize
from os import remove
import numpy as np

# ---------------
# Fixed variables
# ---------------
# Sbox & inverse Sbox
subBytesTable = (
    0x63, 0x7c, 0x77, 0x7b, 0xf2, 0x6b, 0x6f, 0xc5, 0x30, 0x01, 0x67, 0x2b, 0xfe, 0xd7, 0xab, 0x76,
    0xca, 0x82, 0xc9, 0x7d, 0xfa, 0x59, 0x47, 0xf0, 0xad, 0xd4, 0xa2, 0xaf, 0x9c, 0xa4, 0x72, 0xc0,
    0xb7, 0xfd, 0x93, 0x26, 0x36, 0x3f, 0xf7, 0xcc, 0x34, 0xa5, 0xe5, 0xf1, 0x71, 0xd8, 0x31, 0x15,
    0x04, 0xc7, 0x23, 0xc3, 0x18, 0x96, 0x05, 0x9a, 0x07, 0x12, 0x80, 0xe2, 0xeb, 0x27, 0xb2, 0x75,
    0x09, 0x83, 0x2c, 0x1a, 0x1b, 0x6e, 0x5a, 0xa0, 0x52, 0x3b, 0xd6, 0xb3, 0x29, 0xe3, 0x2f, 0x84,
    0x53, 0xd1, 0x00, 0xed, 0x20, 0xfc, 0xb1, 0x5b, 0x6a, 0xcb, 0xbe, 0x39, 0x4a, 0x4c, 0x58, 0xcf,
    0xd0, 0xef, 0xaa, 0xfb, 0x43, 0x4d, 0x33, 0x85, 0x45, 0xf9, 0x02, 0x7f, 0x50, 0x3c, 0x9f, 0xa8,
    0x51, 0xa3, 0x40, 0x8f, 0x92, 0x9d, 0x38, 0xf5, 0xbc, 0xb6, 0xda, 0x21, 0x10, 0xff, 0xf3, 0xd2,
    0xcd, 0x0c, 0x13, 0xec, 0x5f, 0x97, 0x44, 0x17, 0xc4, 0xa7, 0x7e, 0x3d, 0x64, 0x5d, 0x19, 0x73,
    0x60, 0x81, 0x4f, 0xdc, 0x22, 0x2a, 0x90, 0x88, 0x46, 0xee, 0xb8, 0x14, 0xde, 0x5e, 0x0b, 0xdb,
    0xe0, 0x32, 0x3a, 0x0a, 0x49, 0x06, 0x24, 0x5c, 0xc2, 0xd3, 0xac, 0x62, 0x91, 0x95, 0xe4, 0x79,
    0xe7, 0xc8, 0x37, 0x6d, 0x8d, 0xd5, 0x4e, 0xa9, 0x6c, 0x56, 0xf4, 0xea, 0x65, 0x7a, 0xae, 0x08,
    0xba, 0x78, 0x25, 0x2e, 0x1c, 0xa6, 0xb4, 0xc6, 0xe8, 0xdd, 0x74, 0x1f, 0x4b, 0xbd, 0x8b, 0x8a,
    0x70, 0x3e, 0xb5, 0x66, 0x48, 0x03, 0xf6, 0x0e, 0x61, 0x35, 0x57, 0xb9, 0x86, 0xc1, 0x1d, 0x9e,
    0xe1, 0xf8, 0x98, 0x11, 0x69, 0xd9, 0x8e, 0x94, 0x9b, 0x1e, 0x87, 0xe9, 0xce, 0x55, 0x28, 0xdf,
    0x8c, 0xa1, 0x89, 0x0d, 0xbf, 0xe6, 0x42, 0x68, 0x41, 0x99, 0x2d, 0x0f, 0xb0, 0x54, 0xbb, 0x16
    )

invSubBytesTable = (
    0x52, 0x09,	0x6a, 0xd5, 0x30, 0x36,	0xa5, 0x38,	0xbf, 0x40,	0xa3, 0x9e, 0x81, 0xf3, 0xd7, 0xfb,
    0x7c, 0xe3,	0x39, 0x82, 0x9b, 0x2f,	0xff, 0x87,	0x34, 0x8e,	0x43, 0x44, 0xc4, 0xde, 0xe9, 0xcb,
    0x54, 0x7b,	0x94, 0x32, 0xa6, 0xc2,	0x23, 0x3d,	0xee, 0x4c,	0x95, 0x0b, 0x42, 0xfa, 0xc3, 0x4e,
    0x08, 0x2e,	0xa1, 0x66, 0x28, 0xd9,	0x24, 0xb2,	0x76, 0x5b,	0xa2, 0x49, 0x6d, 0x8b, 0xd1, 0x25,
    0x72, 0xf8,	0xf6, 0x64, 0x86, 0x68,	0x98, 0x16,	0xd4, 0xa4,	0x5c, 0xcc, 0x5d, 0x65, 0xb6, 0x92,
    0x6c, 0x70,	0x48, 0x50, 0xfd, 0xed,	0xb9, 0xda,	0x5e, 0x15,	0x46, 0x57, 0xa7, 0x8d, 0x9d, 0x84,
    0x90, 0xd8,	0xab, 0x00, 0x8c, 0xbc,	0xd3, 0x0a,	0xf7, 0xe4,	0x58, 0x05, 0xb8, 0xb3, 0x45, 0x06,
    0xd0, 0x2c,	0x1e, 0x8f, 0xca, 0x3f,	0x0f, 0x02,	0xc1, 0xaf,	0xbd, 0x03, 0x01, 0x13, 0x8a, 0x6b,
    0x3a, 0x91,	0x11, 0x41, 0x4f, 0x67,	0xdc, 0xea,	0x97, 0xf2,	0xcf, 0xce, 0xf0, 0xb4, 0xe6, 0x73,
    0x96, 0xac,	0x74, 0x22, 0xe7, 0xad,	0x35, 0x85,	0xe2, 0xf9,	0x37, 0xe8, 0x1c, 0x75, 0xdf, 0x6e,
    0x47, 0xf1,	0x1a, 0x71, 0x1d, 0x29,	0xc5, 0x89,	0x6f, 0xb7,	0x62, 0x0e, 0xaa, 0x18, 0xbe, 0x1b,
    0xfc, 0x56,	0x3e, 0x4b, 0xc6, 0xd2,	0x79, 0x20,	0x9a, 0xdb,	0xc0, 0xfe, 0x78, 0xcd, 0x5a, 0xf4,
    0x1f, 0xdd,	0xa8, 0x33, 0x88, 0x07,	0xc7, 0x31,	0xb1, 0x12,	0x10, 0x59, 0x27, 0x80, 0xec, 0x5f,
    0x60, 0x51,	0x7f, 0xa9, 0x19, 0xb5,	0x4a, 0x0d,	0x2d, 0xe5,	0x7a, 0x9f, 0x93, 0xc9, 0x9c, 0xef,
    0xa0, 0xe0,	0x3b, 0x4d, 0xae, 0x2a,	0xf5, 0xb0,	0xc8, 0xeb,	0xbb, 0x3c, 0x83, 0x53, 0x99, 0x61,
    0x17, 0x2b,	0x04, 0x7e, 0xba, 0x77,	0xd6, 0x26,	0xe1, 0x69,	0x14, 0x63, 0x55, 0x21, 0x0c, 0x7d
    )

# Round constants
round_constant = (
    0x00000000, 0x01000000, 0x02000000,
    0x04000000, 0x08000000, 0x10000000,
    0x20000000, 0x40000000, 0x80000000,
    0x1B000000, 0x36000000, 0x6C000000,
    0xD8000000, 0xAB000000, 0x4D000000,
    )


# ---------------
# Main action functions
# ---------------
# Xtime
# Used to preform multiplication by x in the Galois field
def xtime(a):
    return (((a << 1) ^ 0x1B) & 0xFF) if (a & 0x80) else (a << 1)


# Byte substitution function
# Substitutes each byte in the state with a byte from the S-Box
def sub_bytes(data, bytesTable):
    for i, row in enumerate(data):
        for j, byte in enumerate(row):
            data[i][j] = bytesTable[byte]
    return data


# Shift rows function
# Shifts the rows of the the matrix to the left.
# Each row is shifted by the number of its index
def shift_rows(array):
    array[:, 1] = np.roll(array[:, 1], -1, axis=0)
    array[:, 2] = np.roll(array[:, 2], -2, axis=0)
    array[:, 3] = np.roll(array[:, 3], -3, axis=0)
    return array


# Inverse shift rows function
# Shifts the rows of the the matrix to the right.
# Each row is shifted by the number of its index
def inv_shift_rows(array):
    array[:, 1] = np.roll(array[:, 1], 1, axis=0)
    array[:, 2] = np.roll(array[:, 2], 2, axis=0)
    array[:, 3] = np.roll(array[:, 3], 3, axis=0)
    return array


# Performs the mix columns layer
def mix_columns(data):
    # mixes a single column
    def mix_single_column(data):
        t = data[0] ^ data[1] ^ data[2] ^ data[3]
        u = data[0]
        data[0] ^= t ^ xtime(data[0] ^ data[1])
        data[1] ^= t ^ xtime(data[1] ^ data[2])
        data[2] ^= t ^ xtime(data[2] ^ data[3])
        data[3] ^= t ^ xtime(data[3] ^ u)

    # mixes all columns using mix_single_column
    def mix(data):
        for i in range(4):
            mix_single_column(data[i])
        return data
    data = mix(data)
    return data


# Preforms the inverse mix columns layer
# This function is similar to the mix_columns function
# but instead preforms the inverse operation.
def inv_mix_columns(data):
    for i in range(4):
        u = xtime(xtime(data[i][0] ^ data[i][2]))
        v = xtime(xtime(data[i][1] ^ data[i][3]))
        data[i][0] ^= u
        data[i][1] ^= v
        data[i][2] ^= u
        data[i][3] ^= v
    mix_columns(data)
    return data


# Adds a padding to ensure a bloke size of 16 bytes
def add_padding(data):
    length = 16 - len(data)
    for i in range(length):
        data.append(0)
    return data, length


# Removes the padding from the data
def remove_padding(data, identifier):
    if identifier[-1] == 0:
        return data
    elif identifier[-1] > 0 and identifier[-1] < 16:
        return data[:-identifier[-1]]
    else:
        raise ValueError('Invalid padding')


# Performs the encryption rounds on the input data matrix
# This function is used for the encryption of data matrixes
# using the expanded keys.
def encryption_rounds(data, round_keys, nr):
    # Inizial add round key
    data = np.bitwise_xor(data, round_keys[0])

    # Rounds 1 to 9 or 1 to 11 or 1 to 13
    # Here each step in one round is performed in a sequence n times
    # where n is the number of rounds minus the last round.
    for i in range(1, (nr - 1)):
        # Sub bytes
        data = sub_bytes(data, subBytesTable)
        # Shift rows
        data = shift_rows(data)
        # Mix columns
        data = mix_columns(data)
        # Add round key
        data = np.bitwise_xor(data, round_keys[i])

    # Final round
    # Identical to the previous rounds, but without mix columns
    data = sub_bytes(data, subBytesTable)
    data = shift_rows(data)
    data = np.bitwise_xor(data, round_keys[nr - 1])

    # Returns the encrypted data
    return data


# Performs the decryption rounds on the input data matrix
# This function is used for the decryption of data matrixes
# using the expanded keys.
def decryption_rounds(data, round_keys, nr):
    # Inizial add round key, inverse shift rows and inverse sub bytes
    data = np.bitwise_xor(data, round_keys[-1])
    data = inv_shift_rows(data)
    data = sub_bytes(data, invSubBytesTable)

    # Rounds 1 to 9 or 1 to 11 or 1 to 13
    # Here each step in one round is performed in a sequence n times
    # where n is the number of rounds minus the last round.
    for i in range(1, (nr - 1)):
        # Add round key
        data = np.bitwise_xor(data, round_keys[-(i+1)])
        # Inverse mix columns
        data = inv_mix_columns(data)
        # Inverse shift rows
        data = inv_shift_rows(data)
        # Inverse sub bytes
        data = sub_bytes(data, invSubBytesTable)

    # Final round
    # Final add round key of final round
    data = np.bitwise_xor(data, round_keys[0])

    # Returns the decrypted data
    return data


# ---------------
# Key expansion setup
# ---------------
# Key expansion function
# This function is used to expand the key to the correct number of round
# keys for the encryption and decryption rounds.
def keyExpansion(key):
    # Format key correctly for the key expansion
    key = [key[i:i+2] for i in range(0, len(key), 2)]

    # Key expansion setup
    # This part determines the number of rounds and the number of words
    # using the key length.
    if len(key) == 16:
        words = key_schedule(key, 4, 11)
        nr = 11
    if len(key) == 24:
        words = key_schedule(key, 6, 13)
        nr = 13
    if len(key) == 32:
        words = key_schedule(key, 8, 15)
        nr = 15

    # Create list for storing the round keys & tmp list for storing
    # for temporary storage.
    round_keys = [None for i in range(nr)]
    tmp = [None for i in range(4)]

    # Formats the words to a list of tuples
    for i in range(nr * 4):
        for index, t in enumerate(words[i]):
            tmp[index] = int(t, 16)  # type: ignore
        words[i] = tuple(tmp)

    # Formats teh words to a list of numpy arrays where each
    # array is a 4x4 matrix representing a round key.
    for i in range(nr):
        round_keys[i] = np.array(words[i * 4] + words[i * 4 + 1] + words[i * 4 + 2] + words[i * 4 + 3]).reshape(4, 4)  # type: ignore

    # Returns the list of round keys and the number of rounds
    return round_keys, nr


# Key schedule (nk = number of colums, nr = number of rounds)
# This function is used to expand the key to the correct number of round
def key_schedule(key, nk, nr):
    # Create list for storing the words and populates the first
    # 4 with the specified key.
    words = [(key[4*i], key[4*i+1], key[4*i+2], key[4*i+3]) for i in range(nk)]

    # Fill out the rest based on previews four words using the fucnitons, rotword,
    # subword and rcon values
    limit = False
    for i in range(nk, (nr * nk)):
        # Get required previous keywords
        temp, word = words[i-1], words[i-nk]

        # If multiple of nk use rot, sub, rcon etc
        if i % nk == 0:
            x = SubWord(RotWord(temp))
            rcon = round_constant[int(i/nk)]
            temp = hexor(x, hex(rcon)[2:])
            limit = False
        elif i % 4 == 0:
            limit = True

        if i % 4 == 0 and limit and nk >= 8:
            temp = SubWord(temp)

        # Xor the two hex values
        xord = hexor(''.join(word), ''.join(temp))
        # Add to list
        words.append((xord[:2], xord[2:4], xord[4:6], xord[6:8]))
    # Return the list of words
    return words


# Takes two hex values and calculates hex1 xor hex2
def hexor(hex1, hex2):
    # Convert to binary
    bin1 = hex2binary(hex1)
    bin2 = hex2binary(hex2)

    # Calculate
    xord = int(bin1, 2) ^ int(bin2, 2)

    # Cut prefix
    hexed = hex(xord)[2:]

    # Leading 0s get cut above, if not length 8 add a leading 0
    if len(hexed) != 8:
        hexed = '0' + hexed

    # Return hex
    return hexed


# Takes a hex value and returns binary
def hex2binary(hex):
    return bin(int(str(hex), 16))


# Takes from 1 to the end, adds on from the start to 1
def RotWord(word):
    return word[1:] + word[:1]


# Selects correct values from sbox based on the current word
# and replaces the word with the new values.
def SubWord(word):
    # Create list for storing the new word
    sWord = []

    # Loop through the current word
    for i in range(4):

        # Check first char, if its a letter(a-f) get corresponding decimal
        # otherwise just take the value and add 1
        if word[i][0].isdigit() is False:
            row = ord(word[i][0]) - 86
        else:
            row = int(word[i][0])+1

        # Repeat above for the seoncd char
        if word[i][1].isdigit() is False:
            col = ord(word[i][1]) - 86
        else:
            col = int(word[i][1])+1

        # Get the index base on row and col (16x16 grid)
        sBoxIndex = (row*16) - (17-col)

        # Get the value from sbox and removes prefix (0x)
        piece = hex(subBytesTable[sBoxIndex])[2:]

        # Check length to ensure leading 0s are not forgotton
        if len(piece) != 2:
            piece = '0' + piece

        # Adds the new value to the list
        sWord.append(piece)

    # Returning word as string
    return ''.join(sWord)


# ---------------
# Running modes setup
# ---------------
# ECB encryption function
def ecb_enc(key, file_path):
    file_size = getsize(file_path)
    round_keys, nr = keyExpansion(key)

    with open(f"{file_path}.enc", 'wb') as output, open(file_path, 'rb') as data:
        for i in range(int(file_size/16)):
            raw = np.array([i for i in data.read(16)]).reshape(4, 4)
            result = bytes((encryption_rounds(raw, round_keys, nr).flatten()).tolist())
            output.write(result)

        if file_size % 16 != 0:
            raw = [i for i in data.read()]  # type: ignore
            raw, length = add_padding(raw)

            result = bytes((encryption_rounds(np.array(raw).reshape(4, 4), round_keys, nr).flatten()).tolist())
            identifier = bytes((encryption_rounds(np.array([0 for i in range(15)] + [length]).reshape(4, 4), round_keys, nr).flatten()).tolist())

            output.write(result + identifier)
        else:
            identifier = bytes((encryption_rounds(np.array([0 for i in range(16)]).reshape(4, 4), round_keys, nr).flatten()).tolist())
            output.write(identifier)
    remove(file_path)


# ECB decryption function
def ecb_dec(key, file_path):
    file_size = getsize(file_path)
    file_name = file_path[:-4]
    round_keys, nr = keyExpansion(key)

    with open(f"{file_name}", 'wb') as output, open(file_path, 'rb') as data:
        for i in range(int(file_size/16) - 2):
            raw = np.array([i for i in data.read(16)]).reshape(4, 4)
            result = bytes((decryption_rounds(raw, round_keys, nr).flatten()).tolist())
            output.write(result)

        data_pice = np.array([i for i in data.read(16)]).reshape(4, 4)
        identifier = np.array([i for i in data.read()]).reshape(4, 4)

        result = (decryption_rounds(data_pice, round_keys, nr).flatten()).tolist()
        identifier = (decryption_rounds(identifier, round_keys, nr).flatten()).tolist()

        result = bytes(remove_padding(result, identifier))

        output.write(result)
    remove(file_path)


# CBC encryption function
def cbc_enc(key, file_path, iv):
    file_size = getsize(file_path)
    vector = np.array([int(iv[i:i+2], 16) for i in range(0, len(iv), 2)]).reshape(4, 4)
    round_keys, nr = keyExpansion(key)

    with open(f"{file_path}.enc", 'wb') as output, open(file_path, 'rb') as data:
        for i in range(int(file_size/16)):
            raw = np.array([i for i in data.read(16)]).reshape(4, 4)
            raw = np.bitwise_xor(raw, vector)
            vector = encryption_rounds(raw, round_keys, nr)
            output.write(bytes((vector.flatten()).tolist()))

        if file_size % 16 != 0:
            raw = [i for i in data.read()]  # type: ignore
            raw, length = add_padding(raw)

            raw = np.bitwise_xor(np.array(raw).reshape(4, 4), vector)
            vector = encryption_rounds(raw, round_keys, nr)

            identifier = np.bitwise_xor(np.array([0 for i in range(15)] + [length]).reshape(4, 4), vector)
            identifier = encryption_rounds(identifier, round_keys, nr)

            output.write(bytes((vector.flatten()).tolist() + (identifier.flatten()).tolist()))
        else:
            identifier = np.bitwise_xor(np.array([0 for i in range(16)]).reshape(4, 4), vector)
            identifier = bytes(((encryption_rounds(identifier, round_keys, nr)).flatten()).tolist())  # type: ignore
            output.write(identifier)  # type: ignore
    remove(file_path)


# CBC decryption function
def cbc_dec(key, file_path, iv):
    iv = np.array([int(iv[i:i+2], 16) for i in range(0, len(iv), 2)]).reshape(4, 4)
    file_size = getsize(file_path)
    file_name = file_path[:-4]
    round_keys, nr = keyExpansion(key)

    with open(f"{file_name}", 'wb') as output, open(file_path, 'rb') as data:
        if int(file_size/16) - 3 >= 0:
            vector = np.array([i for i in data.read(16)]).reshape(4, 4)
            raw = decryption_rounds(vector, round_keys, nr)
            result = np.bitwise_xor(raw, iv)
            output.write(bytes((result.flatten()).tolist()))

            for i in range(int(file_size/16) - 3):
                raw = np.array([i for i in data.read(16)]).reshape(4, 4)
                result = decryption_rounds(raw, round_keys, nr)
                result = np.bitwise_xor(result, vector)
                vector = raw
                output.write(bytes((result.flatten()).tolist()))
        else:
            vector = iv

        data_pice = np.array([i for i in data.read(16)]).reshape(4, 4)
        vector_1, identifier = data_pice, np.array([i for i in data.read()]).reshape(4, 4)

        result = decryption_rounds(data_pice, round_keys, nr)
        identifier = decryption_rounds(identifier, round_keys, nr)

        identifier = np.bitwise_xor(identifier, vector_1)
        data_pice = np.bitwise_xor(result, vector)

        result = bytes(remove_padding((data_pice.flatten()).tolist(), (identifier.flatten()).tolist()))

        output.write(result)
    remove(file_path)


# PCBC encryption function
def pcbc_enc(key, file_path, iv):
    file_size = getsize(file_path)
    vector = np.array([int(iv[i:i+2], 16) for i in range(0, len(iv), 2)]).reshape(4, 4)
    round_keys, nr = keyExpansion(key)

    with open(f"{file_path}.enc", 'wb') as output, open(file_path, 'rb') as data:
        for i in range(int(file_size/16)):
            raw = np.array([i for i in data.read(16)]).reshape(4, 4)
            tmp = np.bitwise_xor(raw, vector)
            vector = encryption_rounds(tmp, round_keys, nr)
            output.write(bytes((vector.flatten()).tolist()))
            vector = np.bitwise_xor(vector, raw)

        if file_size % 16 != 0:
            raw = [i for i in data.read()]  # type: ignore
            raw, length = add_padding(raw)
            raw = np.array(raw).reshape(4, 4)

            tmp = np.bitwise_xor(raw, vector)
            vector1 = encryption_rounds(tmp, round_keys, nr)
            vector = np.bitwise_xor(vector1, raw)

            identifier = np.bitwise_xor(np.array([0 for i in range(15)] + [length]).reshape(4, 4), vector)
            identifier = encryption_rounds(identifier, round_keys, nr)

            output.write(bytes((vector1.flatten()).tolist() + (identifier.flatten()).tolist()))
        else:
            identifier = np.bitwise_xor(np.array([0 for i in range(16)]).reshape(4, 4), vector)
            identifier = bytes((encryption_rounds(identifier, round_keys, nr).flatten()).tolist())  # type: ignore
            output.write(identifier)  # type: ignore
    remove(file_path)


# PCBC decryption function
def pcbc_dec(key, file_path, iv):
    iv = np.array([int(iv[i:i+2], 16) for i in range(0, len(iv), 2)]).reshape(4, 4)
    file_size = getsize(file_path)
    file_name = file_path[:-4]
    round_keys, nr = keyExpansion(key)

    with open(f"{file_name}", 'wb') as output, open(file_path, 'rb') as data:
        if int(file_size/16) - 3 >= 0:
            vector = np.array([i for i in data.read(16)]).reshape(4, 4)
            raw = decryption_rounds(vector, round_keys, nr)
            result = np.bitwise_xor(raw, iv)
            vector = np.bitwise_xor(vector, result)
            output.write(bytes((result.flatten()).tolist()))

            for i in range(int(file_size/16) - 3):
                raw = np.array([i for i in data.read(16)]).reshape(4, 4)
                result = decryption_rounds(raw, round_keys, nr)
                result = np.bitwise_xor(result, vector)
                vector = np.bitwise_xor(raw, result)
                output.write(bytes((result.flatten()).tolist()))
        else:
            vector = iv

        data_pice = np.array([i for i in data.read(16)]).reshape(4, 4)
        vector_1, identifier = data_pice, np.array([i for i in data.read()]).reshape(4, 4)

        result = decryption_rounds(data_pice, round_keys, nr)
        data_pice = np.bitwise_xor(result, vector)

        vector_1 = np.bitwise_xor(vector_1, data_pice)
        identifier = decryption_rounds(identifier, round_keys, nr)
        identifier = np.bitwise_xor(identifier, vector_1)

        result = bytes(remove_padding((data_pice.flatten()).tolist(), (identifier.flatten()).tolist()))

        output.write(result)
    remove(file_path)


# OFB encryption function
def ofb_enc(key, file_path, iv):
    file_size = getsize(file_path)
    round_keys, nr = keyExpansion(key)
    mix = np.array([int(iv[i:i+2], 16) for i in range(0, len(iv), 2)]).reshape(4, 4)
    iv = mix

    with open(f"{file_path}.enc", 'wb') as output, open(file_path, 'rb') as data:
        for i in range(int(file_size/16)):
            raw = np.array([i for i in data.read(16)]).reshape(4, 4)
            mix = encryption_rounds(mix, round_keys, nr)
            result = np.bitwise_xor(raw, mix)
            output.write(bytes((result.flatten()).tolist()))

        if file_size % 16 != 0:
            raw = [i for i in data.read()]  # type: ignore
            raw, length = add_padding(raw)
            raw = np.array(raw).reshape(4, 4)

            if file_size < 16:
                mix = encryption_rounds(iv, round_keys, nr)
            else:
                mix = encryption_rounds(mix, round_keys, nr)
            result = np.bitwise_xor(mix, raw)

            mix = encryption_rounds(mix, round_keys, nr)
            identifier = np.bitwise_xor(np.array([0 for i in range(15)] + [length]).reshape(4, 4), mix)

            output.write(bytes((result.flatten()).tolist() + (identifier.flatten()).tolist()))
        else:
            mix = encryption_rounds(mix, round_keys, nr)
            identifier = np.bitwise_xor(np.array([0 for i in range(16)]).reshape(4, 4), mix)
            output.write(bytes((identifier.flatten()).tolist()))
    remove(file_path)


# OFB decryption function
def ofb_dec(key, file_path, iv):
    iv = np.array([int(iv[i:i+2], 16) for i in range(0, len(iv), 2)]).reshape(4, 4)
    file_size = getsize(file_path)
    file_name = file_path[:-4]
    round_keys, nr = keyExpansion(key)

    with open(f"{file_name}", 'wb') as output, open(file_path, 'rb') as data:
        if int(file_size/16) - 3 >= 0:
            raw = np.array([i for i in data.read(16)]).reshape(4, 4)
            mix = encryption_rounds(iv, round_keys, nr)
            result = np.bitwise_xor(raw, mix)
            output.write(bytes((result.flatten()).tolist()))

            for i in range(int(file_size/16) - 3):
                raw = np.array([i for i in data.read(16)]).reshape(4, 4)
                mix = encryption_rounds(mix, round_keys, nr)
                result = np.bitwise_xor(raw, mix)
                output.write(bytes((result.flatten()).tolist()))
        else:
            mix = iv

        data_pice = np.array([i for i in data.read(16)]).reshape(4, 4)
        identifier = np.array([i for i in data.read()]).reshape(4, 4)

        mix = encryption_rounds(mix, round_keys, nr)
        data_pice = np.bitwise_xor(data_pice, mix)

        mix = encryption_rounds(mix, round_keys, nr)
        identifier = np.bitwise_xor(identifier, mix)

        result = bytes(remove_padding((data_pice.flatten()).tolist(), (identifier.flatten()).tolist()))  # type: ignore

        output.write(result)  # type: ignore
    remove(file_path)

\end{python}

\section{encrypt.py}
\label{app:encrypt.py}

\begin{python}
from PyAES import AES
from sys import argv


# ---------------
# Encryption function
# ---------------
def encrypt(key, file_path, running_mode, iv=None):

    # Input validation
    if (len(key) / 2) not in [16, 24, 32]:
        raise Exception('Key length is not valid')
    elif running_mode in ["CBC", "PCBC", "CFB", "OFB", "CTR", "GCM"]:
        if (len(iv) / 2) != 16 or iv is None:
            raise Exception('IV length is not valid')

    # Running mode selection
    if running_mode == "ECB":
        AES.ecb_enc(key, file_path)
    elif running_mode == "CBC" and iv is not None:
        AES.cbc_enc(key, file_path, iv)
    elif running_mode == "PCBC" and iv is not None:
        AES.pcbc_enc(key, file_path, iv)
    elif running_mode == "OFB" and iv is not None:
        AES.ofb_enc(key, file_path, iv)
    else:
        raise Exception("Running mode not supported")


if __name__ == "__main__":
    encrypt(key=argv[1], file_path=argv[2], running_mode=argv[3], iv=argv[4])

\end{python}

\section{decrypt.py}
\label{app:decrypt.py}

\begin{python}
from PyAES import AES
from sys import argv


# ---------------
# Decryption function
# ---------------
def decrypt(key, file_path, running_mode, iv=None):

    # Input validation
    if file_path[-4:] != ".enc":
        raise Exception('File is not encrypted in known format')
    if (len(key) / 2) not in [16, 24, 32]:
        raise Exception('Key length is not valid')
    elif running_mode in ["CBC", "PCBC", "CFB", "OFB", "CTR", "GCM"]:
        if (len(iv) / 2) != 16 or iv is None:
            raise Exception('IV length is not valid')

    # Running mode selection
    if running_mode == "ECB":
        AES.ecb_dec(key, file_path)
    elif running_mode == "CBC" and iv is not None:
        AES.cbc_dec(key, file_path, iv)
    elif running_mode == "PCBC" and iv is not None:
        AES.pcbc_dec(key, file_path, iv)
    elif running_mode == "OFB" and iv is not None:
        AES.ofb_dec(key, file_path, iv)
    else:
        raise Exception("Running mode not supported")


if __name__ == "__main__":
    decrypt(key=argv[1], file_path=argv[2], running_mode=argv[3], iv=argv[4])

\end{python}

\section{\_\_main\_\_.py}
\label{app:__main__.py}

\begin{python}
from PyAES.encrypt import encrypt
from PyAES.decrypt import decrypt
from getpass import getpass
import PyAES


def main():
    print("-"*85)
    print(r"""      _       ________   ______         _______          _   __
     / \     |_   __  |.' ____ \       |_   __ \        / |_[  |
    / _ \      | |_ \_|| (___ \_| ______ | |__) |_   __`| |-'| |--.   .--.   _ .--.
   / ___ \     |  _| _  _.____`.||______||  ___/[ \ [  ]| |  | .-. |/ .'`\ \[ `.-. |
 _/ /   \ \_  _| |__/ || \____) |       _| |_    \ '/ / | |, | | | || \__. | | | | |
|____| |____||________| \______.'      |_____| [\_:  /  \__/[___]|__]'.__.' [___||__]
                                               \__.'                                 """)
    print("-"*85)
    print(f"Version: {PyAES.__version__}                                      {PyAES.__copyright__}")
    print("-"*85)
    print("""This is a simple AES (Advanced Encryption Standard) implementation in Python-3. It is
a pure Python implementation of AES that is designed to be used as a educational tool
only. It is not intended to be used in any other use case than educational and no
security is guaranteed for data encrypted or decrypted using this tool.""")
    print("-"*85)
    run()


def run():
    action = input("Do you want to encrypt, decrypt or quit? (e/d/q): ")
    if action == "e":
        running_mode = input("Please select cipher running mode (ECB/CBC/PCBC/CFB/OFB/CTR/GCM): ")

        if running_mode == "ECB":
            key = getpass(prompt="Please enter your key: ")
            file_path = input("Please enter path to file: ")
            confirmation = input("Are you sure you want to encrypt this file? (y/n): ")

            if confirmation == "y":
                encrypt(key, file_path, running_mode)
                print("\nEncryption complete!")

            elif confirmation == "n":
                print("Encryption aborted!")
                exit()

            else:
                print("Invalid input!")
                exit()

        elif running_mode in ["CBC", "PCBC", "CFB", "OFB", "CTR", "GCM"]:
            key = getpass(prompt="Please enter your key: ")
            iv = getpass(prompt="Please enter your iv: ")
            file_path = input("Please enter path to file: ")
            confirmation = input("Are you sure you want to encrypt this file? (y/n): ")

            if confirmation == "y":
                encrypt(key, file_path, running_mode, iv)
                print("\nEncryption complete!")

            elif confirmation == "n":
                print("Encryption aborted!")
                exit()

            else:
                print("Invalid input!")
                exit()

        else:
            print("Invalid cipher running mode")
            run()

    elif action == "d":
        running_mode = input("Please select cipher running mode (ECB/CBC/PCBC/CFB/OFB/CTR/GCM): ")

        if running_mode == "ECB":
            key = getpass(prompt="Please enter your key: ")
            file_path = input("Please enter path to file: ")
            confirmation = input("Are you sure you want to decrypt this file? (y/n): ")

            if confirmation == "y":
                decrypt(key, file_path, running_mode)
                print("\nDecryption complete!")

            elif confirmation == "n":
                print("Decryption aborted!")
                exit()

            else:
                print("Invalid input!")
                exit()

        elif running_mode in ["CBC", "PCBC", "CFB", "OFB", "CTR", "GCM"]:
            key = getpass(prompt="Please enter your key: ")
            iv = getpass(prompt="Please enter your iv: ")
            file_path = input("Please enter path to file: ")
            confirmation = input("Are you sure you want to decrypt this file? (y/n): ")

            if confirmation == "y":
                decrypt(key, file_path, running_mode, iv)
                print("\nDecryption complete!")

            elif confirmation == "n":
                print("Decryption aborted!")
                exit()

            else:
                print("Invalid input!")
                exit()

        else:
            print("Invalid cipher running mode")
            run()

    elif action == "q":
        print("Exiting...")
        exit()

    else:
        print("Invalid action (to quit enter 'q')")
        run()


if __name__ == "__main__":
    main()

\end{python}

\chapter{Test kod (Analyze.py)}
\label{app:analyze}

\begin{python}
# ---------------------------------------------------------------
#                   Background structure
# ---------------------------------------------------------------

# Imports the encrypt function from the PyAES module
from PyAES.encrypt import encrypt
# Imports the timer and random number generator
from time import perf_counter
from random import randint
# Imports the remove function for deleting files
from os import remove


# For displaying progress during test runs
def progress_bar(progress, total_progress):
    percent = 100 * (float(progress) / float(total_progress))
    bar_progress = int(100 * (float(progress) / float(total_progress)))

    if bar_progress > 100 or percent > 100:
        bar_progress = 100
        percent = 100

    bar_remaining = 100 - bar_progress
    bar = '#' * bar_progress + '-' * bar_remaining
    print(f"\r[{bar}] {percent:.2f}%", end="\r")
    return progress + 1


# Writes the resulting data to a text file
def write_data(data, name_f):
    # Creates a text file with the specified name
    with open(name_f + '.txt', 'w') as f:
        for i in data:
            # Writes the data to the text file
            f.write(str(i) + '\n')
        f.write(str('\n'))


# Creates a text file with specified size and fills it with random bytes
def setup(count):
        # Creates a text file with the name test_speed.txt
        with open("test_speed.txt", 'wb') as f:
            for j in range(count):
                # Writes random bytes to the text file
                f.write(bytes([randint(0, 255)]))


# Runs the specified function and returns the time it takes to run
def speed_test(count, key, mode, iv):
    # Creates a text file with the specified size
    setup(count)
    # Starts the timer
    start = perf_counter()
    # Executes the function
    encrypt(key, "test_speed.txt", mode, iv=iv)
    # Stops the timer
    end = perf_counter()
    # Deletes the text file
    remove("test_speed.txt.enc")
    return end - start

# Executes the code if the file is run directly
if __name__ == '__main__':
    # Test parameters
    keys = ["2b7e151628aed2a6abf7158809cf4f3c",
            "8e73b0f7da0e6452c810f32b809079e562f8ead2522c6b7b",
    "603deb1015ca71be2b73aef0857d77811f352c073b6108d72d9810a30914dff4"]
    iv = "000102030405060708090a0b0c0d0e0f"
    file_size = 1000000
    runs = 25

    # ---------------------------------------------------------------
    #     Test time difference between 128, 192 and 256 bit keys
    # ---------------------------------------------------------------
    data = []
    progress = 0
    # Loops through the different key sizes
    for i in keys:
        data_tmp = []
        # Runs the test the specified amount of times
        for j in range(runs):
            # Displays the progress
            progress = progress_bar(progress, 150)
            # Runs the test and saves the result
            test = speed_test(file_size, i, 'ECB', iv)
            # Saves the result
            data_tmp.append(test)
        # Writes the the results of every run to a text file
        write_data(data_tmp, 'keys_test_raw ' + str(len(i) * 4))
        # Saves the average of the results
        data.append(sum(data_tmp)/len(data_tmp))
    # Writes the average of the results to a text file
    write_data(data, 'keys_test')

    # ---------------------------------------------------------------
    #      Test time difference between ECB, CBC and OFB modes
    # ---------------------------------------------------------------
    data = []
    data_tmp = []

    # Runs the time test for OFB the specified amount of times
    for i in range(runs):
        # Displays the progress
        progress = progress_bar(progress, 150)
        # Runs the test for OFB
        test = speed_test(file_size, keys[0], 'OFB', iv)
        # Saves the result in a temporary list
        data_tmp.append(test)
    # Writes the the results of every run to a text file
    write_data(data_tmp, 'modes_test_raw OFB')
    # Saves the average of the results in a list and clears the temporary list
    data.append(sum(data_tmp)/len(data_tmp))
    data_tmp = []

    # Runs the time test for CBC the specified amount of times
    for i in range(runs):
        # Displays the progress
        progress = progress_bar(progress, 150)
        # Runs the test for CBC
        test = speed_test(file_size, keys[0], 'CBC', iv)
        # Saves the result in a temporary list
        data_tmp.append(test)
    # Writes the the results of every run to a text file
    write_data(data_tmp, 'modes_test_raw CBC')
    # Saves the average of the results in a list and clears the temporary list
    data.append(sum(data_tmp)/len(data_tmp))
    data_tmp = []

    # Runs the time test for ECB the specified amount of times
    for i in range(runs):
        # Displays the progress
        progress = progress_bar(progress, 150)
        # Runs the test for ECB
        test = speed_test(file_size, keys[0], 'ECB', iv)
        # Saves the result in a temporary list
        data_tmp.append(test)
    # Writes the the results of every run to a text file
    write_data(data_tmp, 'modes_test_raw ECB')
    # Saves the average of the results in a list
    data.append(sum(data_tmp)/len(data_tmp))
    # Writes the averages from every run to a text file
    write_data(data, 'modes_test')

    # Dilsapys the progress finished
    progress = progress_bar(progress, 150)
    # Prints that the tests is completed
    print("\n")
    print("Completed")

\end{python}


% End document
\end{document}